\scanpage{15}
% The first draft, including its heading, is enclosed and cancelled by long strokes.
\deleted{\uline{Cor. 11} Supposons maintenant qu'on ait un syst. projectif \inserted{filtrant strict} $(R_i)_{i\in I}$ de Modules coh. sur $U$, pris affine $U=\operatorname{Spec}A$, et \deleted{supposons} soit $F$ le foncteur sur les $\unsure{\mathcal O_U|\mathscr O_U|A_U}$-Modules}
\deleted{\inserted{$F$ --}}

\uline{Corollaire} 11 \unsure{Soit|Soient} $U,\Omega$ comme dessus, \deleted{$U$ pris affine,} \inserted{Supposons} $S=U=\operatorname{Spec}A$, soit $(P_i)$ \unsure{une|un} \deleted{famille} syst. projectif de $A$-modules, et $F$ le foncteur \unsure{qu'il|qui il} proreprésente. Supposons que l'on ait les deux conditions suivantes

(i) $F(\Omega)$ est un $A$-Module de t.f.

(ii) Pour tt $p\in S$ et tt syst. de paramètres $(f_1,\ldots,f_{\unsure{r|n}})$ de $\deleted{\unsure{\mathcal O|\mathscr O}}A_p$, l'application canonique
\[
 F(\Omega)\otimes A_p/(f_1\unsure{,\ldots,|-}f_r)
 \longrightarrow F(\Omega\otimes A_p/(f_1\unsure{,\ldots,|-}f_r,))
\]
est bijective.

Alors le syst. projectif $(P_i)$ est stationnaire.

De façon précise, (i) \deleted{\unsure{signifie|indique}} \inserted{équivaut à dire} que le système inductif des $\operatorname{Hom}(P_i,\Omega)$ est stationnaire. Soit \deleted{\unsure{$P_i$}} $i_0$ tel que $\operatorname{Hom}(P_{i_0},\unsure{\Omega|L})\overset{\sim}{\longrightarrow}\operatorname{Hom}(P_i,\Omega)$ pour $i\geq i_0$. Alors $P_i\simeq P_{i_0}$ pour $i\geq i_0$.

\uline{Dém} Le diagramme

\diagram{\begin{tikzcd}[ampersand replacement=\&,column sep=small,row sep=large,font=\scriptsize]
\underset{i}{\underset{\longrightarrow}{\lim}}\operatorname{Hom}(P_i,\Omega)\otimes A_p/\underline f A_p \arrow[r,"\sim"] \&
\underset{i}{\underset{\longrightarrow}{\lim}}\operatorname{Hom}(P_i,\Omega_p/\underline f\Omega_p)\\
\operatorname{Hom}(P_{i_0},\Omega)\otimes A_p/\unsure{f|\underline f}A_p \arrow[u,"{\unsure{\wr|\downarrow}}"] \arrow[r] \&
\operatorname{Hom}(P_{i_0},\Omega_p/\underline f\Omega_p) \arrow[u,hook',"\alpha"]
\end{tikzcd}}
% archive sheet number: 14
